% !TeX root = ../../tesis.tex
%%%%%%%%%%%%%%%%%%%%%%%%%%%%%%%%%%%%%%%%%%%%%%%%%%%%%%%%%%%%%%%%%%%
% §6.3.1 — NIVEL DE COBERTURA ESPACIAL (C) EN LOS CASOS CON ANILLO
% (archivo 6-3-2 por la numeración anterior). Redactada 2026-10-22.
%
% Datos (corrida 2026-10-04; RESPUESTAS_PREGUNTAS_ISSUES.md:160,532 y
%   garibelt/comparativo/cobertura_comparativa.csv):
%   C ZMVM 76 = 13.29 % → 13.43 % (0.166 → 0.168); CDMX 53.52 → 53.53 %;
%   municipios 3.92 → 4.09 %; 42 → 41 municipios en cero (Texcoco);
%   7 demarcaciones cambian; Casos 2 y 3 idénticos (mismas estaciones).
%   Figuras usan dominio 141 (0.055 → 0.056), decisión C6-1.
% PENDIENTE #34: recalcular el ámbito CDMX (16 alcaldías) al final.
% Figuras (decisión del autor: las dos en el cuerpo y en el Anexo G
%   a página completa):
%   Clasificacion_Garibelt_Mapa_Cobertura_MB (= _Metro)
%   Comparativo_Garibelt_Cobertura
%%%%%%%%%%%%%%%%%%%%%%%%%%%%%%%%%%%%%%%%%%%%%%%%%%%%%%%%%%%%%%%%%%%
\subsection{Nivel de Cobertura Espacial (\texorpdfstring{$C$}{C})}
\label{subsec:anillo-cobertura}

La cobertura indica qué parte del territorio queda a una distancia caminable
de alguna estación: es el área dentro de \termino{buffers} de 800~m alrededor
de cada parada, como proporción de la superficie de cada demarcación
(§\ref{subsec:indicador-cobertura}). En el Caso~1 la red cubre el 13.29\,\% de
la \zmvm{}, un promedio que reúne situaciones muy distintas. Más de la mitad
de la \cdmx{} (53.52\,\%) está cerca de una estación; en los 60 municipios
conurbados la proporción baja a 3.92\,\% y 42 de ellos no tienen cobertura
(§\ref{subsec:cobertura-espacial}). Normalizado, el indicador vale 0.166 y
queda en el intervalo 0--0.25 del Espectro MJG
(Cuadro~\ref{tab:espectro-mjg-diagnostico}).

\subsubsection{Casos~2 y~3}

En la cobertura los dos casos con anillo dan el mismo resultado. El indicador
depende solo de dónde están las estaciones, no de la velocidad ni del derecho
de vía, y el Anillo MB y el Anillo Metro tienen las mismas 94. Con ellas, la
superficie cubierta crece unos 11~km$^2$ y la cobertura de la \zmvm{} pasa de
13.29 a 13.43\,\%. En la \cdmx{} queda prácticamente igual (53.53\,\%) y en
los municipios conurbados sube a 4.09\,\%. El valor normalizado pasa de 0.166
a 0.168, dentro del mismo intervalo. La
Figura~\ref{fig:anillo-cobertura-mb} muestra la cobertura por demarcación con
el anillo; el mapa del Caso~2 vale también para el Caso~3, porque los dos son
idénticos en este indicador.

\begin{figure}[H]
    \centering
    \includegraphics[width=\textwidth]{%
        Figures/Mapas/Qgis_Resources/RedConAnillo/Escenario_MB/Clasificacion_Garibelt/Clasificacion_Garibelt_Mapa_Cobertura_MB}
    \caption[Cobertura por demarcación en la red con anillo]{%
    Cobertura espacial por demarcación en el Caso~2 (Anillo MB), igual a la
    del Caso~3. El mapa usa las 141 demarcaciones del conjunto de datos; las
    cifras del texto corresponden a las 76 de la \zmvm{}. Versión a página
    completa en la Figura~\ref{fig:anx-g-cobertura-mb}. Corrida: 2026-10-04.}
    \label{fig:anillo-cobertura-mb}
    \fuentefigura{Elaboración propia con \textbf{QGIS};
    datos: \texttt{Apímetro} y \texttt{VFTModel}~\citep{cabrera2025apimetro, cabrera2026VFTModel}.}
\end{figure}

Como se observa en la Figura~\ref{fig:anillo-cobertura-comparativo}, siete
demarcaciones registran algún cambio. Solo en dos se nota: Nezahualcóyotl, que
pasa de 39.8 a 48.6\,\%, y Tlalnepantla de Baz, que pasa de 53.9 a 58.9\,\%;
las dos están en el Estado de México, sobre el norte y el oriente del anillo.
Ecatepec de Morelos, Texcoco, Naucalpan de Juárez, Iztapalapa y Xochimilco
ganan menos de medio punto, y ninguna demarcación cambia de rango en el mapa.
Con el anillo, Texcoco deja la cobertura cero y los municipios sin cobertura
bajan de 42 a 41.

\begin{figure}[H]
    \centering
    \includegraphics[width=\textwidth]{%
        Figures/Mapas/Qgis_Resources/Comparativo_CG/Comparativo_Garibelt_Cobertura}
    \caption[Cobertura por demarcación: red actual y red con anillo]{%
    Cobertura espacial por demarcación en el Caso~1 (A) y en los Casos~2 y~3
    (B), que coinciden porque comparten estaciones. Las demarcaciones
    achuradas son las de mayor ganancia. El mapa usa las 141 demarcaciones del
    conjunto de datos; las cifras del texto corresponden a las 76 de la
    \zmvm{}. Versión a página completa en la
    Figura~\ref{fig:anx-g-cobertura-comparativo}. Corrida: 2026-10-04.}
    \label{fig:anillo-cobertura-comparativo}
    \fuentefigura{Elaboración propia con \textbf{QGIS};
    datos: \texttt{Apímetro} y \texttt{VFTModel}~\citep{cabrera2025apimetro, cabrera2026VFTModel}.}
\end{figure}

\subsubsection{Lectura Comparativa}

El anillo casi no amplía la cobertura, y la razón está en su trazo. El
Periférico recorre la parte de la ciudad que ya concentra más estaciones, así
que en la \cdmx{} los \termino{buffers} de las estaciones nuevas caen casi
siempre sobre área que ya estaba cubierta. Donde el anillo sale a los
municipios, en el norte y el oriente, sí agrega superficie, porque ahí la red
estructurada es escasa. La cobertura mide cercanía, no conexión. Un anillo que
acerca pocas estaciones a territorio nuevo puede, aun así, cambiar los
recorridos dentro de la red, y eso es lo que revisan la Fuerza Capilar, el
tiempo de viaje y la intermediación en los apartados siguientes. Los 41
municipios que siguen sin cobertura están fuera del alcance del \textsc{api-i}
y son el territorio que tendría que atender un anillo exterior
(§\ref{subsec:corredor-anillar-limitaciones}).
